\honest{Mind Projection Fallacy}{Eliezer Yudkowsky}{2008}

On old science fiction magazine covers, alien monsters carry off women in torn dresses. ``Oddly enough, the aliens never go after men in torn shirts.'' Would an alien with a different evolutionary history desire a human woman? ``It seems rather unlikely. To put it mildly.''

The artists did not reason that all minds are alike, so a bug-eyed monster will find women attractive. ``Probably the artist did not even think to ask.'' They thought about the torn dress, not the alien's evolution: tear the dress and she is sexier, and the monster does not enter into it. They saw sexiness as a property of the woman herself, like her height. \nb{The post does not consider that the woman was on the cover to sell the magazine to human buyers, which would explain the cover without any error about minds on the artist's part.}

``Apparently we instinctively represent Sexiness as a direct attribute of the Woman object.'' If your brain uses ``that data structure, or something metaphorically similar to it,'' then sexiness feels like part of the woman, not of the one looking at her, and since she is attractive, the monster will be attracted. Isn't that logical?

E. T. Jaynes, whom I introduce as ``a late grand master of the Bayesian Conspiracy,'' named this error the Mind Projection Fallacy: taking a property of your own mind for a property of the world. Jaynes applied it mainly to probability, which he saw as a state of knowledge and not a feature of objects. ``More about this shortly.'' \nb{The argument is in the next day's post, ``Probability is in the Mind.''}

The error is everywhere: in the dispute over the falling tree, in the magazine cover, in ``Kant's declaration that space by its very nature is flat,'' and in Hume's definition of a priori ideas as those ``discoverable by the mere operation of thought.'' \nb{Kant held that space is a form of the mind's intuition, not a property of things in themselves; what he claimed about the world was that all space we can experience must obey Euclid's geometry. Hume wrote that ``the mind has a great propensity to spread itself on external objects'': he described this very error.}

I end with two stories: a man in a romance with a sentient alien plant who learns it is male, frets, and decides it no longer matters; and, in Foglio and Pollotta's \textsc{Illegal Aliens}, a planet of sentient insects showing a movie poster of a human carrying off a bug in a chiffon dress. ``Just thought I'd mention that.''
